\section{Every large stable standard factor has a nonfaithful complement}
\label{sec:standard-same-quotient-kernel}

\paragraph{Review status and shared provenance.}
D-TD48, reserved question Q2fd, strengthens the maximal-standard
consequence of D-TD43.  The author \texttt{/root/g34\_eightfolds}
proposed using the same-quotient singleton supplied by D-TD45;
root supplied the complete ordinary branching case division and
requested the nonmaximal conclusion.  The argument below reconstructs
all source, quotient, operator and branching interfaces.  The complete
author proof and source-interface reading passes.  Independent
whole-proof review and root acceptance also pass.
No new numerical execution or classification theorem is asserted.

Throughout, a stable tensor factor means that its full matrix
algebra is normalized by the original finite coefficient image.
All complementary factors are retained.

\begin{theorem}[D-TD48]
Let \(C\to\mathbf P^1\) be an actual connected finite \(J\)-Galois
cover, with faithful absolutely irreducible primitive complex
coefficient representation \(V\).  Suppose there is an actual
surjection and actual finite linear factorization
\[
 x:J\twoheadrightarrow T,\qquad T\in\{A_m,S_m\},\qquad
 V=(\operatorname{Std}_m\circ x)\otimes U,\qquad
 m\ge2^{30}.
 \tag{SQK.1}
\]
Write \(b=\dim U\ge1\), \(e=(m-1)b\), and suppose the actual
coefficient inertia tuple has
\[
                    \sum_i\operatorname{rank}(1-V(g_i))=2e+6.
 \tag{SQK.2}
\]
Then \(U\) is not faithful.  In fact
\[
  1\ne K_0:=\ker U,\qquad
  x|_{K_0}\text{ is injective},\qquad
  x(K_0)\supseteq A_m.
 \tag{SQK.3}
\]
Thus \(K_0\) acts transitively on the natural \(m\) letters.
The result requires neither maximality of \(m\), an original
coefficient-field hypothesis, nor a support floor or degree bound
for the entire complement.

If, in addition, \(b\ge2\), \(V\) has a model over its exact
imaginary quadratic character field K, and the actual sixfold
variation has signature \((3,3)\) and period rank nine, the
stronger transitive-kernel interface of HD-SS01 and HD-SS06
supplies actual K-models for \(U\) and its restriction to a
natural letter stabilizer, with the same full image \(U(J)\).
Write \(\mu_U=\dim\operatorname{Hom}_J(U,H^1(C,\mathbf C))\).
If this source has smallest coefficient degree in the stated
exact-K, rank-nine, natural signed Jacobian-determinant class,
then
\[
                              \mu_U>0.
 \tag{SQK.4}
\]
When the displayed standard factor is a singleton extraction
layer in one fixed complete original D-TD12 decomposition,
D-TD44 consequently applies to its entire remaining partner.
These latter conclusions retain their additional field, period,
polarization, decomposition and minimality hypotheses.
\end{theorem}

\subsection*{The ordinary same-quotient tensor obstruction}

We first give the needed elementary representation-theoretic
argument, including its one-box branching input.

\begin{lemma}
For every \(m\ge3\) and partition \(\lambda\vdash m\) other than
\((m)\) and \((1^m)\), the ordinary \(S_m\)-module
\[
                         \operatorname{Std}_m\otimes S^\lambda
 \tag{SQK.5}
\]
is reducible.  In particular, if \(W\) is a nontrivial ordinary
irreducible \(A_m\)-module extending to \(S_m\), one can choose an
extension \(\widetilde W\) for which
\(\operatorname{Std}_m\otimes\widetilde W\) has a nonzero proper
\(S_m\)-invariant subspace.
\end{lemma}

\begin{proof}
We use the ordinary irreducible parametrization and Frobenius
character formula with the following exact primary locators:
Fulton--Harris, Formula~4.10, printed p.~49, Proposition~4.37
and its complete proof, pp.~56--57, and the one-box induction
and branching statements (4.42)--(4.43), pp.~58--59.
\cite{HK:FultonHarris1991}
Here is a direct reconstruction of the one-box step, so no
solution of an exercise or general Littlewood--Richardson rule
is assumed.

For a partition \(\mu\), use at least \(m+1\) variables, pad
\(\mu\) by zeros, and write
\[
 a_{\mu+\delta}(X)=\det(X_j^{\mu_i+r-i}),\qquad
 \delta=(r-1,r-2,\ldots,0),\qquad
 s_\mu=a_{\mu+\delta}/a_\delta .
 \tag{SQK.6}
\]
This is the Schur alternant definition in Appendix~A, (A.4),
printed p.~454.  Expanding the determinant into signed monomials
and multiplying each monomial by \(p_1=\sum_jX_j\) gives
\[
       p_1a_{\mu+\delta}
          =\sum_{i=1}^r a_{\mu+\delta+\epsilon_i}.
 \tag{SQK.7}
\]
The exponents in \(\mu+\delta\) strictly decrease.  Increasing
one exponent by one either produces equality with its predecessor,
in which case that determinant is zero, or gives precisely a
partition obtained by adding one box to \(\mu\).  There is no
reordering sign and each allowed diagram occurs once.  Hence
\(p_1s_\mu=\sum_{\nu\succ\mu}s_\nu\), where \(\nu\succ\mu\)
means addition of one box.

For clarity the induction identity is equally elementary.
For a class type \(\rho\), put
\(z_\rho=\prod_j j^{a_j}a_j!\), where \(a_j\) counts its
j-cycles.  The Frobenius formula identifies the irreducible
characteristic \(\sum_\rho\chi^\mu(\rho)p_\rho/z_\rho\)
with \(s_\mu\).  The induced character from \(S_{m-1}\) at
an element of \(S_m\) is its number of fixed letters times
the character on the remaining letters.  Since
\(z_\rho=a_1z_{\rho\setminus1}\) when \(a_1>0\), its
characteristic is \(p_1s_\mu\).  Thus (SQK.7), and then
Frobenius reciprocity, give both multiplicity-free rules
\[
 \operatorname{Ind}_{S_{m-1}}^{S_m}S^\mu
       =\bigoplus_{\nu\succ\mu}S^\nu,\qquad
 \operatorname{Res}_{S_{m-1}}^{S_m}S^\lambda
       =\bigoplus_{\mu\prec\lambda}S^\mu .
 \tag{SQK.8}
\]

The natural permutation representation is
\(\mathbf1\oplus\operatorname{Std}_m\).  The tensor identity,
or its character evaluated at each fixed natural coset, gives
\[
  (\mathbf1\oplus\operatorname{Std}_m)\otimes S^\lambda
       \simeq
       \operatorname{Ind}_{S_{m-1}}^{S_m}
            \operatorname{Res}_{S_{m-1}}^{S_m}S^\lambda.
 \tag{SQK.9}
\]
Let \(c\) be the number of removable corners of \(\lambda\).
By (SQK.8) and reciprocity, the multiplicity of \(S^\lambda\)
on the right is \(c\), and its multiplicity in (SQK.5) is
\(c-1\).  If \(c\ge3\), there are at least two copies, so
the tensor product is reducible.  If \(c=2\), it contains
one copy of \(S^\lambda\), and its dimension
\((m-1)\dim S^\lambda\) is strictly larger than
\(\dim S^\lambda\).  Semisimplicity supplies a nonzero
complement, again proving reducibility.

It remains that \(c=1\), so \(\lambda=(a^d)\) is a rectangle.
The omitted row and column cases are exactly \(d=1\) and
\(a=1\).  For \(a,d\ge2\), removing its unique corner gives
\(\mu=(a^{d-1},a-1)\).  Its three addable boxes produce
\(\lambda\) and the two distinct partitions
\[
       \alpha=(a+1,a^{d-2},a-1),\qquad
       \beta=(a^{d-1},a-1,1).
 \tag{SQK.10}
\]
Therefore (SQK.8)--(SQK.9) give exactly
\(\operatorname{Std}_m\otimes S^\lambda
 \simeq S^\alpha\oplus S^\beta\).
Both summands are nonzero.  This completes every case.

For the final assertion, an ordinary extendible nontrivial
\(A_m\)-module is the irreducible restriction of some
\(S^\lambda\), with \(\lambda\ne\lambda'\).  Neither row
nor column restricts nontrivially.  Apply the proved result
to that extension.  This assertion does not apply to an
individual split constituent of a self-conjugate shape;
that case will instead be excluded by HD-FT03.
\end{proof}

\subsection*{The actual projective quotient and common natural map}

\begin{proof}[Proof of D-TD48]
Assume for contradiction that \(U\) is faithful.
If \(b=1\), the finite group J embeds in \(\mathbf C^*\)
and is abelian, contrary to its surjection onto T.
Hence \(b\ge2\).

Irreducibility of V implies irreducibility of U.  A nontrivial
direct-sum system of imprimitivity for U would tensor with
the entire standard factor to give such a system for V.
Thus the finite projective image
\[
                         P_U=U(J)/\bigl(U(J)\cap\mathbf C^*\bigr)
 \tag{SQK.11}
\]
is primitive in its represented action.
If \(U(g)\) is scalar, faithfulness of U implies that g
is central in J.  The group T has trivial center, so
\(x(g)=1\).  Consequently x factors as an actual surjection
\(P_U\twoheadrightarrow T\).

Apply the coordinate consequence of accepted D-TD45 to
this stable whole primitive factor U in the original V.
Its proof uses the complete D-TD12 decomposition of the
actual image \(P_U\), with all its coarse matrix algebras
embedded in the original U-algebra.  The external standard
factor is retained.  At \(m\ge2^{30}\), the Heisenberg
coordinate is excluded by D-TD43, and a killed nonabelian
core would give an actual tensor-permutation quotient,
excluded by D-TD45 without a restriction on its orbit size.
Thus the conclusion is an actual singleton coarse factor W,
an entire residual factor R, and a commutative quotient
diagram
\[
 U=W\otimes R\quad\hbox{up to scalar choices},\qquad
 J\longrightarrow P_U\longrightarrow P_W
       \xrightarrow[\sim]{\alpha}T,\qquad
 x=\alpha\circ\operatorname{pr}_W\circ\operatorname{pr}_U .
 \tag{SQK.12}
\]
The algebra of W is individually stable.  Its represented
finite perfect central core E covers \(A_m\) and acts
irreducibly and nontrivially.  No normal lift of E inside
the original J is required.  In particular (SQK.12) is the
same original natural map x, not an unrelated isomorphic
alternating quotient.

\paragraph{The two exceptional module types.}
For \(m\ge8\), the accepted covering-group classification
leaves the ordinary \(A_m\) modules and genuine spin modules
of \(2.A_m\).  The exact covering locator is
Hoffman--Humphreys, Chapter~4 opening, printed p.~31,
as used in HD-FT20 and HD-ML01.\cite{BI:HHM}
In the spin case, W satisfies HD-FT28: its algebra is
stable, its projective image contains and normalizes the
represented \(A_m\), and the whole external partner is
\(\operatorname{Std}_m\otimes R\).  Since \(m\ge2^{30}>10^6\),
that theorem contradicts the original sphere budget
(SQK.2).  It uses no maximality or bound on R.

An ordinary irreducible \(A_m\)-module is either an
extendible Specht restriction or a split constituent.
The exact alternatives and their compatible-normalizer
scopes are those of HD-FT03.  A split constituent has
only scalar-times-inner compatible operators: an odd
automorphism interchanges the two inequivalent constituents,
and a projective scalar discrepancy would be a character
of the perfect group \(A_m\) and hence trivial.  Thus if W
is split, (SQK.12) supplies precisely the actual stable
\(A_m\)-quotient interface of HD-FT03, with the whole
partner \(\operatorname{Std}_m\otimes R\) and all scalar
phases retained.  Its cutoff is \(m>3500000\), which
follows from \(m\ge2^{30}\).  This also contradicts
(SQK.2).  In particular a split module is not erroneously
put through the ordinary tensor-reducibility lemma.

\paragraph{An ordinary extendible module and its actual operators.}
The remaining W is a nontrivial ordinary irreducible
\(A_m\)-module.  Use \(\alpha\) in (SQK.12) to label its
core by that same \(A_m\), and choose the \(S_m\)-extension
\(\sigma\) supplied by the preceding lemma.
For \(g\in J\), let \(A_g\) be any actual local linear
representative of its projective action on W.
It normalizes the perfect represented core honestly:
the derived subgroup of \(\mathbf C^*E\) is E.
Its action on the ordinary core, under the fixed
identification, is conjugation by \(x(g)\).
Thus \(A_g\) and \(\sigma(x(g))\) implement the same
automorphism of E.  Perfectness kills a possible scalar
discrepancy on the core, and Schur's lemma gives
\[
                         A_g=c_g\,\sigma(x(g)),
                         \qquad c_g\in\mathbf C^* .
 \tag{SQK.13}
\]
This argument applies to the odd coset too if \(T=S_m\);
it is not a choice of unrelated natural labels at
different operators.

After absorbing \(c_g\) in the entire remaining factor,
the original coefficient operator has the exact form
\[
 V(g)=
  \bigl(\operatorname{Std}_m(x(g))\otimes\sigma(x(g))\bigr)
                         \otimes B_g .
 \tag{SQK.14}
\]
There is no need to assert that \(g\mapsto B_g\) is an
ordinary representation.  The preceding lemma provides
a fixed nonzero proper \(S_m\)-invariant subspace L in
\(\operatorname{Std}_m\otimes W\).
Equation (SQK.14) makes \(L\otimes R\) a nonzero proper
J-invariant subspace of V, for every scalar phase and
every residual factor.  This contradicts the assumed
irreducibility of V.  All possibilities under faithful U
are exhausted, so \(K_0=\ker U\ne1\).

\paragraph{Kernel and source consequences.}
If an element of \(K_0\) also lies in \(\ker x\), it
acts trivially on V and is the identity.  Thus
\(x|_{K_0}\) is injective.  Its nontrivial normal image
in \(A_m\) or \(S_m\) contains \(A_m\), proving
(SQK.3) and transitivity.  These conclusions are purely
group-theoretic consequences inside the actual source;
they make no coefficient-field or polarization assertion.

Under the separately stated exact-K hypotheses, this
actual normal kernel is precisely the stronger subgroup
interface in HD-SS01/HD-SS06.  It supplies the same-field
models and equality of the full partner image on a
natural letter stabilizer.  If \(b\ge2\) and \(\mu_U=0\),
the post-vanishing part of HD-SS06 gives an actual
degree-b rational-base exact-K source with the same
rank-nine variation and natural signed determinant.
Its fixed coefficient-line and quotient-pullback comparison
scales the rank-six Hermitian form by a positive rational
number c, so the determinant scale is
\(c^6=\operatorname{Nm}_{K/\mathbf Q}(c^3)\).
This contradicts the specified minimality and proves
(SQK.4).  The same complete-decomposition hypotheses then
permit D-TD44; no such decomposition or minimality is
needed for the first nonfaithfulness assertion.
\end{proof}

\paragraph{Normalization and retained boundaries.}
If a stable standard factor was initially specified up to
scalars, use the fixed actual representation
\(\operatorname{Std}_m\circ x\) and absorb its scalar
phases into U.  The tensor identity for V and the
multiplicativity of this fixed standard factor make
the resulting U an actual representation of the finite J.
Its kernel is the kernel used in this theorem.  A sign
twist of the chosen standard extension is absorbed in
the same way; it is not silently imposed on the partner.

The theorem removes the maximality and \(N_*\) restrictions
from the earlier faithful-complement exclusion.  It does
not give an upper bound on \(m\) or \(b\), classify another
extraction component, eliminate positive partner cohomology,
or assign (SQK.2) to the partner alone.  Other signatures
and the remaining nonfaithful positive-partner sources
remain separate.  A scalar change of tensor representatives
is never treated as a change of the original inertia tuple,
field or natural polarization.
